\makeatletter
\@ifundefined{keepfig}{\newenvironment{keepfig}{\par\edef\keepfig@pi{\the\parindent}\noindent
  \begin{minipage}{\linewidth}\parindent=\keepfig@pi\relax}{\end{minipage}\par}}{}
\makeatother
\section{Общий многочлен степени $n$}

Когда мы говорим, что корни многочлена $aX^2+bX+c$ равны
\[
\frac{-b\pm\sqrt{b^2-4ac}}{2a},
\]
мы думаем об $a,b,c$ как о \emph{символах}: при любых конкретных значениях $a,b,c$ формула даёт корни конкретного
уравнения. В этом разделе мы докажем, что для корней <<общего многочлена>> степени $\ge5$ подобной формулы в радикалах нет.

\begin{definition}[5.39b]
Пусть $t_1,\dots,t_n$~--- независимые переменные (символы). \emph{Общим многочленом степени $n$} над полем $F$
называется
\[
f(X)=X^n-t_1X^{n-1}+t_2X^{n-2}-\dots+(-1)^nt_n\in F[t_1,\dots,t_n][X].
\]
Мы рассматриваем его как многочлен от $X$ с коэффициентами в поле $F(t_1,\dots,t_n)$.
\end{definition}

Знаки выбраны так, чтобы $t_r$ играл роль $r$-й элементарной симметрической функции корней (см.\ \rf{5.37}(c)). Покажем,
что группа Галуа $f$ над $F(t_1,\dots,t_n)$ равна $S_n$; тогда теорема \rf{5.34} даст сформулированное выше утверждение (по
крайней мере в характеристике $0$).

\begin{theorem}[5.40]
Группа Галуа общего многочлена степени $n$ равна $S_n$. Точнее, $f$ сепарабелен над $F(t_1,\dots,t_n)$, и группа Галуа
его поля разложения действует на $n$ корнях $f$ всеми возможными перестановками.
\end{theorem}

\emph{Идея.} В \rf{5.39} мы уже знаем многочлен с группой $S_n$: это
\[
g(X)=X^n-p_1X^{n-1}+\dots+(-1)^np_n=\prod_{i=1}^n(X-X_i)\in F(p_1,\dots,p_n)[X].
\]
Он получается из $f$ подстановкой $t_i\mapsto p_i$. Нужно убедиться, что эта подстановка не <<склеивает>> ничего лишнего,
т.\,е.\ задаёт \emph{изоморфизм} полей коэффициентов; тогда многочлены $f$ и $g$ устроены одинаково.

\begin{proof}
\emph{Шаг 1: гомоморфизм $\varphi\colon F[t_1,\dots,t_n]\to F[p_1,\dots,p_n]$, $t_i\mapsto p_i$, инъективен.} Пусть
$P=\sum_dc_d\,t_1^{d_1}\cdots t_n^{d_n}\ne0$. По формуле \eqref{eq:s11-leadmon} старший моном произведения
$p_1^{d_1}\cdots p_n^{d_n}$ равен $X_1^{d_1+\dots+d_n}X_2^{d_2+\dots+d_n}\cdots X_n^{d_n}$, и для разных наборов
$d=(d_1,\dots,d_n)$ эти мономы разные. Поэтому для различных мономов $m_1(t)\ne m_2(t)$ старшие мономы многочленов
$m_1(p_1,\dots,p_n)$ и $m_2(p_1,\dots,p_n)$ различны, и сокращений быть не может: наибольший из старших мономов
слагаемых $c_d\,p_1^{d_1}\cdots p_n^{d_n}$ (с $c_d\ne0$) входит в $P(p_1,\dots,p_n)$ с коэффициентом $c_d\ne0$ (подробно
это разобрано в \rf{5.37}(b)). Значит, $\varphi(P)=P(p_1,\dots,p_n)\ne0$.

\emph{Шаг 2: изоморфизм полей частных.} Гомоморфизм $\varphi$ сюръективен по определению: кольцо $F[p_1,\dots,p_n]$
порождено над $F$ элементами $p_i$. По шагу 1 он~--- изоморфизм колец, и поэтому продолжается до изоморфизма полей частных
\[
\Phi\colon F(t_1,\dots,t_n)\xrightarrow{\ \sim\ }F(p_1,\dots,p_n),\qquad \Phi(a/b)=\varphi(a)/\varphi(b).
\]
Именно здесь нужна инъективность: при $b\ne0$ должно быть $\varphi(b)\ne0$, иначе делить нельзя. (Корректность:
если $a/b=a'/b'$, то $ab'=a'b$, откуда $\varphi(a)\varphi(b')=\varphi(a')\varphi(b)$.) Если бы ядро $\varphi$ было ненулевым,
то никакого продолжения не было бы: ненулевой элемент ядра обратим в $F(t_1,\dots,t_n)$, а перейти должен в $0$.

\emph{Шаг 3: $\Phi$ переводит $f$ в $g$.} Применяя $\Phi$ к коэффициентам, получаем
$X^n-t_1X^{n-1}+\dots+(-1)^nt_n\mapsto X^n-p_1X^{n-1}+\dots+(-1)^np_n=g(X)$.

\emph{Шаг 4: перенос поля разложения и группы Галуа.} Будем рассматривать $F(X_1,\dots,X_n)$ как расширение поля
$F(t_1,\dots,t_n)$ посредством вложения $\Phi\colon F(t_1,\dots,t_n)\to F(p_1,\dots,p_n)\subset F(X_1,\dots,X_n)$. В нём
многочлен $f$ (т.\,е.\ $\Phi f=g$) раскладывается на линейные множители $X-X_i$, и его корни $X_1,\dots,X_n$ порождают
поле. Значит, $F(X_1,\dots,X_n)$~--- поле разложения $f$ над $F(t_1,\dots,t_n)$, и любое поле разложения $L$ многочлена $f$
над $F(t_1,\dots,t_n)$ изоморфно ему над $F(t_1,\dots,t_n)$ (\rf{2.13}); такой изоморфизм $\tilde\Phi\colon
L\to F(X_1,\dots,X_n)$ продолжает $\Phi$ и переводит корни $f$ в корни $g$. Отображение
$\sigma\mapsto\tilde\Phi\circ\sigma\circ\tilde\Phi^{-1}$~--- изоморфизм $\Gal(L/F(t_1,\dots,t_n))\to
\Gal(F(X_1,\dots,X_n)/F(p_1,\dots,p_n))$, согласованный с действием на корнях. По \rf{5.39} последняя группа~--- это $S_n$,
действующая всеми перестановками корней $X_1,\dots,X_n$. Наконец, корни $X_1,\dots,X_n$ попарно различны, так что $f$
сепарабелен.
\end{proof}

\begin{keepfig}
\[
\begin{tikzcd}[column sep=huge,row sep=large]
L\arrow[r,cnew,"\tilde\Phi","\sim"'] & F(X_1,\dots,X_n)\\
F(t_1,\dots,t_n)\arrow[u,hook]\arrow[r,cold,"\Phi","\sim"'] & F(p_1,\dots,p_n)\arrow[u,hook,"S_n"']
\end{tikzcd}
\]
\figcap{Изоморфизм $\Phi$ ($t_i\mapsto p_i$) переводит $f$ в $g$; он продолжается до изоморфизма $\tilde\Phi$ полей разложения
(красная стрелка), и группа Галуа $f$ отождествляется с группой $S_n$ перестановок переменных.}
\end{keepfig}

\begin{remark}[5.40a]
(a) Инъективность $\varphi$ означает, что $p_1,\dots,p_n$ \emph{алгебраически независимы} над $F$: между ними нет
никаких полиномиальных соотношений. Это можно увидеть и иначе: поле $F(X_1,\dots,X_n)$ алгебраично (даже конечно) над
$F(p_1,\dots,p_n)$, поэтому степень трансцендентности $F(p_1,\dots,p_n)$ над $F$ равна $n$, и $n$ образующих
$p_1,\dots,p_n$ обязаны быть алгебраически независимы (см.\ главу 8).

(b) Если подставить вместо $t_i$ элементы, между которыми есть соотношения (например, конкретные числа), то
инъективность пропадает, и группа Галуа может уменьшиться. Например, при $t_1=-1$, $t_2=1$, $t_3=-1$, $t_4=1$
получается круговой многочлен $X^4+X^3+X^2+X+1=\Phi_5$ с группой $(\Z/5\Z)^\times\cong\Z/4\Z$ (см.\ упражнение 5-1).
Однако для <<большинства>> рациональных значений $t_i$ группа остаётся равной $S_n$ (это следует из теоремы Гильберта о
неприводимости, см.\ \rf{5.41}).
\end{remark}

\begin{corollary}[5.40b]
Пусть $\operatorname{char}F=0$ и $n\ge5$. Тогда общий многочлен степени $n$ не разрешим в радикалах над
$F(t_1,\dots,t_n)$.
\end{corollary}

\begin{proof}
По теореме \rf{5.40} его группа Галуа~--- $S_n$, а $S_n$ при $n\ge5$ неразрешима (её единственная нетривиальная собственная
нормальная подгруппа~--- простая неабелева группа $A_n$). Остаётся применить теорему Галуа \rf{5.34}.
\end{proof}

Смысл следствия: не существует <<формулы>>, выражающей корни через коэффициенты $t_1,\dots,t_n$ с помощью
арифметических операций и извлечения корней. Такая формула означала бы, что корни $f$ лежат в радикальном расширении
поля $F(t_1,\dots,t_n)$, т.\,е.\ что $f$ разрешим в радикалах. При $n\le4$ группа $S_n$ разрешима, и формулы
действительно существуют (формулы Кардано для $n=3$ и Феррари для $n=4$, см.\ \rf{5.35d}).

Теорема \rf{5.40} верна над любым полем: общий многочлен сепарабелен и имеет группу Галуа $S_n$. Поэтому при
$\operatorname{char}F=p>0$ по \rf{5.35b} общий многочлен степени $n\ge5$ не разрешим и в обобщённых радикалах (\rf{5.35a}).

\begin{aside}[5.41: обратная задача теории Галуа]
Группа $S_n$ встречается как группа Галуа над $\Q$ (см.\ \rf{4.32}), а каждая конечная группа $G$ вкладывается в
некоторую $S_n$ (теорема Кэли). Если $E/\Q$~--- расширение Галуа с группой $S_n$ и $G\subset S_n$, то $E/E^G$~---
расширение Галуа с группой $G$, а $E^G$~--- конечное расширение $\Q$. Итак, \emph{каждая конечная группа является группой
Галуа над некоторым конечным расширением $\Q$}. Но является ли каждая конечная группа группой Галуа \emph{над самим
$\Q$}, т.\,е.\ группой Галуа некоторого $f\in\Q[X]$? Это \emph{обратная задача теории Галуа}; она до сих пор не решена.

Программа Гильберта--Нётер для её решения была такой. Гильберт доказал: если $G$~--- группа Галуа некоторого расширения
$E\supset\Q(t_1,\dots,t_n)$ ($t_i$~--- символы), то $G$ бесконечно много раз встречается как группа Галуа над $\Q$. Для
доказательства $E$ реализуют как поле разложения многочлена $f(X)\in\Q[t_1,\dots,t_n][X]$ и показывают, что при
бесконечно многих значениях $t_i$ получающийся многочлен из $\Q[X]$ имеет группу Галуа $G$. Эмми Нётер предположила
следующее: если $G\subset S_n$ действует на $F(X_1,\dots,X_n)$ перестановками $X_i$, то
$F(X_1,\dots,X_n)^G\cong F(t_1,\dots,t_n)$ для некоторых символов $t_i$ (для $G=S_n$ это верно: по теореме \rf{5.38}
$F(X_1,\dots,X_n)^{S_n}=F(p_1,\dots,p_n)$, а $p_1,\dots,p_n$ алгебраически независимы по \rf{5.37}(b)). Тогда по
теореме Гильберта все конечные группы встречались бы над $\Q$. Однако Суон в 1969 году доказал, что над $\Q$ гипотеза
неверна для циклической группы порядка $47$. Поэтому этот подход не может доказать, что все конечные группы встречаются
как группы Галуа над $\Q$, хотя других подходов это не исключает.
\end{aside}

\begin{aside}[5.42: геометрическая картина]
Возьмём $F=\CC$ и рассмотрим подмножество $S\subset\CC^{n+1}$, заданное уравнением
\[
X^n-T_1X^{n-1}+\dots+(-1)^nT_n=0 .
\]
Это комплексное многообразие размерности $n$: уравнение однозначно выражает $T_n$ через $(X,T_1,\dots,T_{n-1})$, так что
$S$~--- график многочлена от этих переменных и $S\cong\CC^n$. Рассмотрим проекцию
\[
\pi\colon S\to\CC^n,\qquad (x,t_1,\dots,t_n)\mapsto(t_1,\dots,t_n).
\]
Её слой над точкой $(a_1,\dots,a_n)$~--- множество корней многочлена $X^n-a_1X^{n-1}+\dots+(-1)^na_n$. Дискриминант
$D(f)$ многочлена $f(X)=X^n-T_1X^{n-1}+\dots+(-1)^nT_n$~--- многочлен из $\CC[T_1,\dots,T_n]$ (см.\ \rf{5.37}(c)); пусть
$\Delta\subset\CC^n$~--- множество его нулей. Над каждой точкой $\CC^n\setminus\Delta$ лежит ровно $n$ точек $S$ (корни
простые), и $S\setminus\pi^{-1}(\Delta)\to\CC^n\setminus\Delta$~--- накрытие: по теореме о неявной функции простые
корни локально голоморфно зависят от коэффициентов. Обход петель в $\CC^n\setminus\Delta$ переставляет корни; группа
всех получающихся перестановок (группа монодромии накрытия) равна $S_n$~--- топологический аналог теоремы \rf{5.40}.
\end{aside}

\paragraph{Исторический комментарий.} Ещё во II тысячелетии до н.\,э.\ вавилоняне умели решать квадратные
уравнения. Формулы для корней кубического уравнения (дель Ферро, Тарталья, Кардано) и уравнения четвёртой степени
(Феррари, который ввёл кубическую резольвенту и воспользовался результатом Кардано) были найдены в XVI веке. Следующие
почти три столетия продолжались безуспешные попытки найти аналогичные формулы для уравнений высших степеней, пока Руффини
(1799, с пробелами в доказательстве) и Абель (1824) не доказали, что их не существует.
